\section{Discussion, Limitations and Conclusion}

\noindent\textbf{Practical recipe.}
For a new tactile sensor and known object meshes, first select a small set of registered contacts spanning the intended contact states. Reserve a separate labeled validation set and report its size alongside the training budget. Train a teacher and a residual aligner on the available training labels; use the aligner to assign mesh-derived depth, normal and pose targets to additional real images. Train a student on the labeled and aligned frames. On our objects, one alignment round without confidence filtering is sufficient; repeated rounds give little further improvement.

The renderer-free pipeline is a useful starting point: it retains the within-session depth gains and avoids appearance-model training. When the deployment session differs from the training session, our results support adding synthetic data from a renderer retrained on the labeled and aligned real pairs. At $K=100$ on sensor A, this lowers second-session contact MAE from $0.288$ to $0.252\,$mm. If real contacts of a new object are unavailable but its mesh is known, synthetic pose coverage is also useful: the held-out-object study reduces rotation error from $64.7^\circ$ to $5$--$7^\circ$. Optional test-time alignment can further refine pose when object identity and mesh are available, but its larger residuals across sessions should be considered for the intended application. These recommendations concern the evaluated ridge-pattern contacts; they do not establish a universal label threshold for other sensors or geometries.

\noindent\textbf{Reporting.}
Report the real data used by every component, including the renderer, teacher, aligner and checkpoint-selection procedure. Distinguish labeled training pairs, unlabeled images, validation labels and images inspected during manual calibration. A renderer fitted to all available training pairs can exceed a downstream $K$-shot budget without seeing a validation image. In leave-object-out evaluation, access to real images of the held-out objects is a separate protocol violation. Both should be distinguished from access to the held-out meshes, which is permitted in our geometry-transfer study.

Under this accounting, the all-pairs renderer's $24$--$26\%$ depth-MSE reduction over $K$-shot real-only training is a result at a larger label budget, not a gain from simulation at a cost of $K$ labels. Report metrics that distinguish depth magnitude from contact localization, since all-pixel MSE can improve while contact MAE or IoU worsens. Also identify training duration, seed aggregation, checkpoint selection and test-time mesh use. Contact-region rendering error can help screen appearance models, but its observed correlation with downstream error pools multiple budgets and renderer variants; it should complement downstream evaluation rather than replace it.

\noindent\textbf{Limitations.} MeshTac requires the identity and mesh of each object during pseudo-labeling and optional test-time refinement. Our evaluation covers 12 ridge-pattern objects and does not establish performance on curved or textured contacts, unknown objects, or a manipulation task. GT is mesh-derived on both sensors; sensor-B targets are re-registered public labels. Validation labels are also needed for checkpoint selection. Alignment can fail, particularly across sessions, and repeated ridge patterns leave some translations ambiguous. Sensor-A object identity and session appearance are confounded because each object has one training capture session. The physical renderer uses a geometric approximation to gel deformation; its manual contact-parameter tuning also needs to be included in the data-access account.

\noindent\textbf{Conclusion.} The results indicate that mesh-pose accuracy is a key bottleneck in turning unlabeled real contacts into supervision: refining the pose improves both label quality and student prediction, and the depth gains persist without synthetic images. With $10$--$25$ training labels per object, MeshTac reduces contact MAE by $13$--$20\%$ and $33$--$36\%$ on the two sensors. At the reported operating points, it matches or improves on real-only training with twice as many labels. Accounting for renderer supervision is essential to interpreting these gains. Synthetic data adds robustness across sessions, while the same aligner provides optional test-time pose refinement. Future work will test less constrained contact geometry and downstream manipulation.
